%%
%% Berliner Hochschule für Technik --  Abschlussarbeit
%%
%% Kapitel 4
%%
%%

\chapter{Technische Anforderungen an Endgeräte}
Die CueLens-App wird für das Betriebssystem Android entwickelt. Als Endgeräte für die Teilnahme an der Studie kommen Smartphones und Tablets infrage. Die Anforderungen unterscheiden sich danach, ob nur der aktuelle Studienprototyp mit lokal gebündelten Reizbildern und festgelegten Labelpaaren oder zusätzlich eine KI-gestützte Erkennung selbst aufgenommener Alltagsbilder ausgeführt wird.

\section{Grundfunktionen des Studienprototyps}
Die grundlegende Funktionalität der App stellt moderate Anforderungen an die Hardware der Endgeräte. Bilddarstellungen im Vollbild, ein Startbildschirm mit Countdown, Buttons mit Grafik oder Text, eine Slider-Eingabe für Rauchverlangen und kurze HTTP-POST-Übertragungen sind auf gängigen Android-Smartphones möglich. Auch die typischen Displaygrößen gängiger Smartphones sind für den vorgesehenen Zweck akzeptabel. Android-Tablets sind gleich gut oder besser geeignet als Smartphones. Die App fixiert die Darstellung auf Hochformat, sodass Cue-Bild, Antwortoptionen und Craving-Slider in einer kontrollierten Anordnung präsentiert werden können.

Die Downloadgröße des Android-Packages (APK) liegt unter 50 MB, wenn nur die Grundfunktionen der App und die lokal gebündelten Studienressourcen enthalten sind. Der Payload aller 20 Übertragungen von Selbstberichten beträgt in der aktuellen prototypischen Form insgesamt weniger als 200 KB. Für eine auswertbare Studienversion kommen noch ein pseudonymer Teilnahmetoken, ein Durchgangsindex, die experimentelle Bedingung und gegebenenfalls technische Zählerstände hinzu. Auch dann bleibt die zu übertragende Datenmenge gering, weil keine Bilder, Audiodaten oder Sensordaten übertragen werden.

Der lokale Persistenzbedarf ist ebenfalls gering. Die App speichert den Studienfortschritt, den nächsten erlaubten Startzeitpunkt und die zufällig erzeugte Reihenfolge der Cue-Matching-Aufgaben in \texttt{SharedPreferences}. Diese Datenstruktur besteht aus wenigen Ganzzahlen, Zeitstempeln und einer kommagetrennten Indexliste. Sie ist robust gegenüber App-Neustarts und benötigt keine lokale Datenbank. Aus Datenschutzsicht ist relevant, dass die automatische Android-Sicherung deaktiviert ist, damit lokale Fortschrittsdaten nicht über allgemeine Gerätesicherungen in weitere Kontexte übertragen werden \cite{bsi_tr03161_mobile_2024}.

Die Rechenlast des Studienprototyps entsteht vor allem durch Bilddekodierung, Skalierung und Compose-Rendering. Da alle Stimuli als App-Ressourcen ausgeliefert und die Cue-Labeling-Wortpaare fest codiert sind, ist keine On-Device-Inferenz erforderlich. Die prototypische Operationalisierung mit statischen Ressourcen ist daher alltagstauglicher, energieärmer und auf mehr Endgeräten lauffähig als eine Variante, die in jeder Studiensituation ein neuronales Bildmodell ausführt. Diese Beschränkung ist methodisch zu dokumentieren, weil die App dadurch keine individuellen Alltags-Cues erkennt, sondern einen festen Stimuluspool präsentiert.

\section{KI-Funktionalität}
Das KI-Feature stellt deutlich höhere Ansprüche an Rechen- und Speicherkapazität. Vortrainierte Classifier-Modelle können mehrere hundert MB groß sein und würden die App-Downloadgröße entsprechend erhöhen. Die Klassifizierung eines Bildes dauerte in einem Test mit dem Modell MobileCLIP S2 ungefähr fünf Sekunden auf einem Preis-Leistungs-Smartphone. Der zusätzliche Speicherbedarf lag in der Größenordnung mehrerer hundert MB; die praktische Präzision blieb unter den Erwartungen. Ein eigens für diesen Zweck trainiertes und für mobile Inferenz optimiertes Modell könnte Modellgröße, Speicherbedarf, Laufzeit und Erkennungsqualität verbessern.

In Bezug auf die Android-Version sind die Anforderungen des KI-Features moderat, solange etablierte mobile Inferenzbibliotheken verwendet werden. Für LiteRT wird mindestens das SDK-Level 23 benötigt \cite{google_litert_android_2026}. Diese SDK-Version wurde bereits 2015 veröffentlicht. Die reine Studienfassung ohne On-Device-KI ist von dieser zusätzlichen Modellanforderung unabhängig und benötigt nur die üblichen Android- und Netzwerkfunktionen einer nativen App.
